\scanpage{347}
% top centre
\mspage{263}
% lines 1–8
\uline{Corollaire} 4.10.8 (du \unsure{Th.|th.} 4.10.1) Soit $N\longrightarrow J$ un foncteur cofibrant, tel que les fibres soient des catégories accessibles (4.8), les foncteurs \unsure{coch.} de base $\lambda_*:N_{j'}\longrightarrow N_j$ \unsure{accessible|accessibles}, enfin \deleted{les} $J$ équivalente à une petite catégorie. Alors la catégorie $M=\underline{\unsure{\Gamma|\mathscr T}}(N/J)$ est accessible.
% lines 8–14
Plus précisément, soit $\pi_0$ un cardinal \inserted{infini} tel que \deleted{l'on ait}\inserted{les} \deleted{$\unsure{\mathrm{card}\,J\leq}$}\inserted{catégories} $N_j$ soient $\pi_0$-accessibles, et \inserted{que} les foncteurs $\lambda_*$ soient $\pi_0$-accessibles, et soit $\pi$ un cardinal \inserted{infini} tel que l'on ait \textcircled{a} $\pi^{\pi_0}=\pi$, \inserted{\textcircled{b} $\forall\lambda:j'\longrightarrow j$ dans $J$, $\lambda_*:\unsure{\mathcal E|\mathscr E|\mathcal C}_j\unsure{\mathrel{-}|\longrightarrow}\mathcal E_{j'}$ applique $(\mathcal E_{\unsure{j|j'}})_{\pi_0}$ dans $(\mathcal E_j)_{\unsure{\pi|\pi_0}}$} \deleted{\unsure{et}}\textcircled{\unsure{c|b}} $\mathrm{card}\,\unsure{\mathrm{réd}|\mathrm{red}}\,J\leq\pi$,
% lines 15–16
Alors \deleted{\illegible{1}} $M$ est $\pi$-accessible, et plus précisément, on a
% line 17
\[
 M\unsure{\approxeq|\simeq}\operatorname{Ind}_{\pi}(M_\pi),
\]
% lines 18–20
où $M_\pi$ est la sous-catégorie \unsure{pleine} de $M$ formée des sections \inserted{$F$ de $N$ sur $J$} telles que l'on ait $F(j)\in(\mathcal E_{\deleted{\unsure{\pi}}\unsure{j|\pi}})_\pi\quad\forall j\in\operatorname{Ob}J$,
% archive sheet number: 346
